\section{Simultaneous annuli at prescribed volume scales}
\label{ann:section}

A local estimate with probability close to one must work along every
macroscopic route.  Its volume label and its Euclidean radius are
different parameters.  This section matches them using circle averages
and obtains a linear number of successful, distinct radii around every
point.  All selections of labels are made before revealing the fresh
fields in the selected shells.

Write \(A(a,b)=\{z:a<|z|<b\}\), and fix three shells
\begin{equation}
 E=\overline{A(1/8,4)},\qquad E^+=A(1/16,8),\qquad
 U=A(1/32,16).
 \label{ann:shells}
\end{equation}
The readouts below depend on field restrictions and local kernels in
\(E\); \(U\) is the fresh-field buffer.  Any fixed finite collection of
strictly buffered shell geometries can be placed inside shells with
these roles by changing their fixed radii.  Constants are then allowed
to depend on that collection.

The reference law consists of a zero-boundary real GFF on \(B(0,64)\)
and an independent whole-plane imaginary GFF modulo its angular
period, restricted to \(E\), with the local kernels of
Theorem~\ref{ker:local-kernels}.  The logarithmic covariance
normalization is
\(\Cov(h(z),h(w))=\log|z-w|^{-1}+\text{smooth}\).
The Dirichlet inner product is
\((f,g)_\nabla=(2\pi)^{-1}\int\nabla f\cdot\nabla g\,\dd z\).
An imaginary representative has a circle-average phase in a fixed
interval of length equal to the period.  Subtracting this circle
average leaves the usual centered Gaussian differences.

\subsection{Uniform comparison on a strict interior}

\begin{lemma}[Conditional transfer]\label{ann:conditional-transfer}
Let \(h^0,\widetilde h^0\) be independent zero-boundary GFFs on \(U\).
Add functions \(g,\widetilde g\) harmonic on \(E^+\), with
\[
 \|g\|_{L^2(E^+)}+\|\widetilde g\|_{L^2(E^+)}\le H,
\]
and bounded real and imaginary constants.  On \(E\), retain any fixed
family of the same local output kernels under both field laws.
There is a function \(\omega_H:[0,1]\to[0,1]\), with
\(\omega_H(\varepsilon)\to0\) as \(\varepsilon\downarrow0\), such that
every event of reference probability at most \(\varepsilon\) has
probability at most \(\omega_H(\varepsilon)\) under this shifted fresh
law.  Conversely, for each \(\eta>0\) there is \(p(H,\eta)>0\)
such that every event of reference probability at least \(\eta\) has
shifted fresh probability at least \(p(H,\eta)\).
The constants are uniform in the volume label and in the measurable
event.  They depend on a bound for the two added constants.
\end{lemma}

\begin{proof}
Choose a smooth cutoff supported in \(E^+\) and equal to one near
\(E\).  For a harmonic function \(g\), its product with this cutoff
is an extension in \(H^1_0(U)\), equal to \(g\) on \(E\), whose
Dirichlet norm is at most \(C\|g\|_{L^2(E^+)}\).
Indeed, multiply \(\Delta g=0\) by \(\chi^2g\), integrate by parts,
and use \(2ab\le a^2/2+2b^2\) to bound
\(\int\chi^2|\nabla g|^2\) by
\(4\int g^2|\nabla\chi|^2\); use nested cutoffs for the product
extension.  Constants have extensions of bounded norm by the same
cutoff.  Apply this to both fields in the product Dirichlet space.

If \(P_f\) is the fresh law shifted by such an extension \(f\),
and \(P_0\) the unshifted fresh law, Cameron--Martin and
Cauchy--Schwarz give, for every event \(A\),
\begin{equation}
 P_f(A)\le e^{\|f\|_\nabla^2/2}P_0(A)^{1/2},
 \qquad
 P_f(A)\ge e^{-\|f\|_\nabla^2}P_0(A)^2.
 \label{ann:cm-two-directions}
\end{equation}
These inequalities still hold after adjoining the same conditional
kernel: the likelihood ratio is a function of the fields alone.

Decompose the reference pair on \(U\) by the domain Markov property.
Its harmonic parts, after separating the bounded imaginary phase,
have almost surely finite \(L^2(E^+)\) norms.  Their extensions are
independent of the two fresh parts.  Choose a bound \(H_0\) whose
extension event has probability at least \(1/2\).  Integrating the
second inequality in \eqref{ann:cm-two-directions} gives
\[
 P_{\rm ref}(A)\ge\tfrac12e^{-H_0^2}P_0(A)^2.
\]
The first inequality then bounds the target probability by a constant
times \(P_{\rm ref}(A)^{1/4}\).

For the other direction, if \(P_{\rm ref}(A)\ge\eta\), choose
\(H_\eta\) so that the reference extension exceeds \(H_\eta\)
with probability at most \(\eta/2\).  Integrating the first
inequality gives
\[
 P_0(A)\ge e^{-H_\eta^2}(\eta/2)^2.
\]
The second inequality gives the asserted positive lower bound,
for example \(e^{-H_t^2-2H_\eta^2}(\eta/2)^4\), where \(H_t\)
bounds the target extension norm.  All bounds concern the field
laws and therefore are independent of the output label.
\end{proof}

There is one finite-order issue in applying this lemma.  The
unclocked imaginary input may allow several local traversal orders,
whereas a marginal test uses the actual order.

\begin{lemma}[All feasible orders]\label{ann:all-orders}
For the local arrays of Lemma~\ref{ker:finite-order}, events with
actual-order failure probability at most \(\varepsilon\), uniformly
in their labels, admit augmented all-order events with failure
probability at most \(\omega(\varepsilon)\), where
\(\omega(\varepsilon)\to0\).  Success of the augmented event implies
success for the actual order.  Conditional on both whole fields,
the augmented arrays in disjoint buffered regions have product laws.
\end{lemma}

\begin{proof}
Let \(J\) be the local unclocked imaginary input.  For each feasible
order \(o\) let \(w_o(J)>0\) be its conditional probability and
\(K_o\) its local output kernel.  Lemma~\ref{ker:finite-order}
identifies these weights and kernels: the weights and finite feasible
set depend on \(J\), independently of the real field and volume
label.  Adjoin the other order coordinates as in that lemma, retaining
the actual coordinate, so that the conditional vector law is
\(\prod_oK_o\).

For any \(\zeta>0\), choose \(J_0<\infty\) and \(\delta>0\)
so that, with reference probability at least \(1-\zeta\), there
are at most \(J_0\) feasible orders and all their weights are at
least \(\delta\). This is an exhaustion of the full traversal list for
the fixed protection, under one imaginary marginal law common to all labels.
Neither the bound nor the minimum weight is asserted uniformly over
all imaginary inputs.
On this event, the conditional probability that some coordinate
fails is at most
\[
 \sum_oK_o(\text{failure at }o)
 \le\delta^{-1}\sum_ow_o(J)K_o(\text{failure at }o).
\]
After integration, the all-order failure probability is at most
\(\zeta+\varepsilon/\delta\).  Choose \(\zeta\) first and then
\(\varepsilon\).  The product assertion is the augmented product
assertion of Lemma~\ref{ker:finite-order}.  The same argument works
for a finite joint array of readouts in one buffer.
\end{proof}

\subsection{A Gaussian budget along displaced scale chains}

Fix \(\epsilon_0=2^{-13}\) and \(s=\log 2^{16}\).
At radius \(r_j=e^{-t_j}\), \(t_j=t_0+js\), use the square grid
\(\epsilon_0r_j\Z^2\).
Given a finest-grid point \(y\), let \(c_j(y)\) be its nearest
point of this grid, resolving ties deterministically.
Thus \(|c_j(y)-y|\le\epsilon_0r_j\).
Set
\begin{equation}
 U_j=c_j+r_jU,\qquad
 H_j=h_{32r_j}(c_j),
 \label{ann:coarse-values}
\end{equation}
where the subscript denotes circle average.  These fixed choices
ensure that the closures of \(U_j\) are disjoint: each smaller buffer,
and its circle of radius \(32r_j\), is strictly inside the hole of
every larger buffer.  The circle at the larger scale is outside its
own buffer and all the smaller ones.

Work first with a zero-boundary GFF in a disk \(D\), and with
centers in a fixed compact subset of \(D\).  For sufficiently small
largest radius, all displayed circles are inside \(D\).
Let \(\mathcal F\) contain both fields outside \(\bigcup_jU_j\).
The domain Markov decompositions are
\[
 h=h_j^0+\mathcal H_j,\qquad
 \widetilde h=\widetilde h_j^0+\widetilde{\mathcal H}_j
 \quad\hbox{on }U_j.
\]
The fresh fields are mutually independent and independent of
\(\mathcal F\); the harmonic functions and \(H_j\) are
\(\mathcal F\)-measurable.  For an imaginary lift define
\(\widetilde H_j\) by the same outer circle.  Put
\[
 G_j(u)=\mathcal H_j(c_j+r_ju)-H_j,\qquad
 \widetilde G_j(u)=
   \widetilde{\mathcal H}_j(c_j+r_ju)-\widetilde H_j,
 \quad u\in E^+.
\]
The imaginary differences are independent of the choice of lift.

\begin{lemma}[Square budget]\label{ann:gaussian-budget}
There are constants \(C,\Lambda\), depending only on the fixed
shell geometry and \(s\), such that the centered Gaussian vector
\[
 Z=\bigl((H_j-H_{j-1})_{1\le j\le m},
              (G_j,\widetilde G_j)_{0\le j\le m}\bigr)
 \quad\hbox{in}\quad
 \R^m\oplus\bigoplus_{j=0}^m L^2(E^+)^2
\]
has covariance operator norm at most \(\Lambda\) and trace at most
\(C(m+1)\), once the largest radius is sufficiently small.
The threshold for that radius may depend on the ambient compact set
and disk.  Consequently, when \(m\le C_0N\), for each \(R<\infty\)
one can choose \(A<\infty\) so that
\begin{equation}
 \PP\left[
  \sum_{j=1}^m(H_j-H_{j-1})^2+
  \sum_{j=0}^m\bigl(\|G_j\|_2^2+\|\widetilde G_j\|_2^2\bigr)
       >AN\right]\le e^{-RN}.
 \label{ann:square-tail}
\end{equation}
The bound is uniform in the finest-grid point \(y\).
\end{lemma}

\begin{proof}
We verify the covariance assertion, including the effect of moving
centers.  Evaluation of \(\mathcal H_j\) at \(c_j+r_ju\)
is the boundary pairing against harmonic measure
\(\nu_{j,u}\) on \(\partial U_j\).  For \(u\in E^+\), its
density on the two boundary circles, after scaling, is bounded
uniformly, as are its derivatives on any slightly smaller fixed
shell containing \(E\).  The charge representing \(G_j(u)\)
is \(\nu_{j,u}-\sigma_j\), where \(\sigma_j\) is uniform measure
on the circle of radius \(32r_j\).  It has total mass zero,
bounded total variation, and support at distances comparable to
\(r_j\) from \(y\).  The increment \(H_j-H_{j-1}\) is represented
by \(\sigma_j-\sigma_{j-1}\), with the same properties at these
two adjacent scales.  Smooth circle densities make their logarithmic
self-energies finite and uniformly bounded after scaling.  The
additive \(\log r_j^{-1}\) cancels because the charges have mass
zero.  The ratio of the adjacent radii is fixed.

For two nonadjacent scales, the smaller charge has support of
diameter \(O(r_k)\), at distance comparable to \(r_j\) from
the larger charge.  Subtract the logarithmic kernel at one point
of the smaller support.  Its gradient there is \(O(r_j^{-1})\);
the total-mass cancellation gives a covariance bound
\(C r_k/r_j\).  For increments, associate the charge to its
larger of the two radii.  Enlarging the constant to cover adjacent
scales, every covariance block therefore has norm at most
\[
 C e^{-s(|j-k|-2)_+}.
\]
For the \(L^2\) blocks, the same estimate follows from the uniform
pointwise covariance bound and the fixed area of \(E^+\).
For the smooth part of the Green function, cancellation in both
variables bounds the block by \(C_Dr_jr_k\), with the larger
adjacent radius used for an increment.  Its operator norm is at
most \(C_D\sum_jr_j^2\); it is uniformly bounded once the largest
radius is small.  Pin terms for the imaginary whole-plane field
are sums of a function of either variable and disappear against
the zero-mass charges.

The row sums of the exponential block bound are uniformly finite.
Cauchy--Schwarz, or the block Schur bound, gives the asserted operator
norm.  Integrating the diagonal covariance bound gives the trace
estimate.  This also constructs the profiles as Gaussian
\(L^2\)-valued variables.

If \(\lambda_\ell\) are the covariance eigenvalues, for
\(0<t<(4\Lambda)^{-1}\),
\[
 \EE e^{t\|Z\|^2}
 =\prod_\ell(1-2t\lambda_\ell)^{-1/2}
 \le \exp\!\left(2t\sum_\ell\lambda_\ell\right)
 \le e^{C'(m+1)}.
\]
Finite-dimensional approximation and monotone convergence justify
the formula.  Markov's inequality, followed by increasing \(A\),
proves \eqref{ann:square-tail}.
\end{proof}

\subsection{Alignment and simultaneous coverage}

A unit test at label \(v>0\), transplanted to center \(c\) and
radius \(r\), is evaluated with real field
\begin{equation}
 h^{c,r,v}(u)=h(c+ru)+Q\log r-\gamma^{-1}\log v.
 \label{ann:normalized-field}
\end{equation}
Its imaginary input is pulled back by the same affine map, with its
phase modulo the period.  Its local output is the corresponding
array readout under the affine and relative-volume identities of
Proposition~\ref{ker:field-transfer}.  Thus its deterministic cost
unit is part of the test.  This definition allows arbitrary
deterministic extremal-length units.

\begin{theorem}[Aligned annuli]\label{ann:aligned-annuli}
Fix \(\theta>0\) and \(M>0\).  For each integer \(N\) in any
unbounded deterministic set, let
\[
 I_N\subset[N,2N]\cap\Z,\qquad |I_N|\ge\theta N,\qquad
 v_i=M4^{-i}.
\]
Consider local unit events \(\mathcal G_{N,i}\), \(i\in I_N\),
for the arrays of Theorem~\ref{ker:local-kernels}, on the fixed
shell \(E\). Their reference probabilities are evaluated under the
single normalized ordinary field law specified above, at every label.
The events may vary with \(N,i\); the local field law and full
buffer geometry do not.  There are constants
\[
 0<a<B<\infty,\quad c>0,\quad \varepsilon_*>0
\]
depending only on \(\gamma,\theta\) and the fixed buffer geometries,
with the following implication.  If every \(\mathcal G_{N,i}\)
has reference probability at least \(1-\varepsilon_*\), then in
an ordinary field chart every compact set \(K\) has, almost surely
for all sufficiently large admissible \(N\), a collection of good
transplanted tests such that:
\begin{enumerate}
\item There are at most \(CN\) radii, all in
      \([e^{-BN},e^{-aN}]\).  At each radius the centers belong
      to a deterministic square grid of spacing
      \(\epsilon_0r\). At each fixed radius the buffers have overlap
      bounded by a fixed constant.
\item Every \(x\in K\) has at least \(cN\) good tests at
      distinct radii \(r\), with centers \(c_r\) satisfying
      \(|x-c_r|<r/256\).  In particular \(x\) lies with strict
      clearance in each central hole \(B(c_r,r/128)\).
\item Each test has a label in \(I_N\), is evaluated by
      \eqref{ann:normalized-field}, and succeeds for the actual
      traversal order.  At most a fixed number of labels are
      used at any one center and radius; one may retain a single
      successful test there.
\end{enumerate}
Under the ordinary Gaussian law, failure probabilities are bounded
by \(C_Ke^{-c_KN}\) and hence are summable.  The almost sure
eventual conclusion persists under any fixed locally absolutely
continuous change of the joint field law, with the same local
kernels.  It holds simultaneously over countably many compact
sets and test families satisfying the stated hypotheses.
\end{theorem}

\begin{proof}
There are two different grid counts in the proof. We control coarse
circle averages at the first and last radii on their own grids.
The harmonic-profile and fresh-shell failures are then made to have
arbitrarily large exponential rates, so that they can be summed over
the much finer grid indexing every displaced scale chain.

Choose \(2<\beta<Q\) and set
\[
 a=\frac{\log4}{4\gamma(Q+\beta)},\qquad
 B=\frac{4\log4}{\gamma(Q-\beta)},\qquad
 t_j=aN+js,\quad
 m=\left\lfloor\frac{(B-a)N}{s}\right\rfloor .
\]
All buffers are contained in a fixed compact neighborhood of \(K\)
for large \(N\).  Use the grids and nearest-center chains above.
The number of finest-grid points needed to cover \(K\) is at most
\(C_Ke^{2BN}\).

First control the two endpoints on their own grids.
The circle averages there have variance \(t_j+O_K(1)\).
There are at most \(C_Ke^{2t_j}\) centers on the grid at scale
\(j\).  The Gaussian tail bound therefore gives, for \(j=0,m\),
\begin{equation}
 \PP[\text{some such center has }|H_j|>\beta t_j]
 \le C_K e^{-(\beta^2/2-2)t_j}
 \le C_K e^{-c_KN}.
 \label{ann:endpoint-tail}
\end{equation}
Increasing \(C_K\) absorbs the bounded variance correction.
On their complement, for every chain the exterior-measurable
process
\[
 Y_j=\gamma(Qt_j-H_j)+\log M
\]
satisfies \(Y_0<N\log4\) and \(Y_m>2N\log4\) for sufficiently
large \(N\).  The fixed multiplier \(M\) changes only this last
threshold.

Choose \(R>2B+3\), and then the budget \(A\) in
Lemma~\ref{ann:gaussian-budget}.  On its good event,
\begin{equation}
 \sum_{j=1}^m(Y_j-Y_{j-1})^2\le A_1N
 \label{ann:jump-budget}
\end{equation}
for a deterministic \(A_1\), since \(t_j-t_{j-1}=s\)
and \(m=O(N)\).

Assign each \(i\in I_N\) to the first \(j\) for which
\(Y_j\ge i\log4\).  Discard the assignment if
\(Y_j-i\log4>L\), where \(L\ge\log4\) will be fixed shortly.
For one upward jump \(\Delta=Y_j-Y_{j-1}\), the number of
discarded lattice levels is zero unless \(\Delta>L\), and in
that case is at most
\[
 \Delta/\log4+1\le 2\Delta/\log4
                 \le\frac{2\Delta^2}{L\log4}.
\]
The sum in \eqref{ann:jump-budget} therefore bounds the number
discarded by \(2A_1N/(L\log4)\).  Choose \(L\) so large
that this is at most \(\theta N/8\).
Every remaining assignment satisfies
\begin{equation}
 |Y_j-i\log4|\le L,
 \label{ann:alignment}
\end{equation}
and a single \(j\) has at most
\(K_L=\lceil2L/\log4\rceil+2\) such assignments.

Now choose \(H\) so large that
\(K_LA/H^2\le\theta/8\).
There are at most \(AN/H^2\) indices at which
\(\|G_j\|_2^2+\|\widetilde G_j\|_2^2>H^2\).
Deleting their labels leaves at least \(\theta N/2\) labels
at at least
\[
 \kappa_0N,\qquad \kappa_0=\theta/(2K_L),
\]
distinct radii.  Select one label at each such radius, for example
the smallest.  Endpoint tests, budget tests, label assignments, and
this selection are all \(\mathcal F\)-measurable.  This is why
the coarse circles in \eqref{ann:coarse-values} were placed outside
every fresh annulus.

For a selected pair \((j,i)\), the conditional real field in
\eqref{ann:normalized-field} is a fresh field on \(U\) plus
\[
 G_j+\left(H_j-Qt_j-\gamma^{-1}\log M
                                  +i\gamma^{-1}\log4\right).
\]
The constant in parentheses has absolute value at most \(L/\gamma\).
The imaginary conditional field is a fresh field plus
\(\widetilde G_j\) and a phase in a bounded interval.
Lemma~\ref{ann:all-orders} first makes the reference all-order
failure probability small.  Lemma~\ref{ann:conditional-transfer}
then makes its conditional target failure probability at most any
prescribed \(p_*>0\), by choosing \(\varepsilon_*\) small enough.
These choices are made after
\(A,L,H\) and the geometric constants have been fixed.

The selected augmented tests are conditionally independent:
integrate the fresh-field product law against the conditional
product of local kernels in the disjoint buffers.  The random
selection causes no change, since it is \(\mathcal F\)-measurable
and takes values in a finite deterministic list.
With \(\ell\ge\kappa_0N\) selected radii, a union bound over
subsets shows
\begin{equation}
 \PP[\text{fewer than }\ell/2\text{ tests succeed}\mid\mathcal F]
 \le 2^\ell p_*^{\ell/2}
 \le 2^{m+1}p_*^{\kappa_0N/2}.
 \label{ann:bernoulli-tail}
\end{equation}
Choose \(p_*\) to make this at most \(e^{-RN}\).
Together with the square-budget bound this has an arbitrarily
large exponential rate.  Taking the union over the at most
\(C_Ke^{2BN}\) finest-grid chains leaves a summable error.
The endpoint errors are added using \eqref{ann:endpoint-tail};
they were union-bounded on their own grids and do not incur the
finest-grid entropy.

For any \(x\in K\), choose a finest-grid point \(y\) within
\(\epsilon_0r_m\).  At every scale of its chain,
\[
 |x-c_j(y)|\le\epsilon_0r_m+\epsilon_0r_j
                  \le 2\epsilon_0r_j<r_j/256.
\]
Its at least \(\kappa_0N/2\) successes give the required distinct
radii.  Take the union of the successful tests produced by all
chains.  At a fixed center and radius, the alignment interval
allows at most \(K_L\) labels.  If several succeed, retaining any
one preserves coverage, because their central holes are identical.
The radius-\(r\) buffers lie in disks of radius \(16r\) about
grid points spaced by \(\epsilon_0r\), so their overlap is at
most a constant depending only on \(\epsilon_0\).

Borel--Cantelli proves eventual coverage.  This event belongs to
the restriction of the fields and kernels to any fixed neighborhood
of \(K\); finitely many large initial radii can be discarded.
A fixed absolutely continuous field-law change preserves this
probability-one event by Proposition~\ref{ker:field-transfer}.
It need not preserve the displayed summable numerical probability
bound.  Finally intersect the probability-one events for a countable
compact exhaustion and countably many prescribed test families.
\end{proof}

For transfer back to a jointly convergent discrete array, fix \(N\)
and its finite collection of centers, labels and strict buffered tests
first. On the event that these tests succeed, their strict inequalities
and protected incidences hold in the approximating surfaces for all
sufficiently large discrete indices, by
\Cref{ker:local-kernels,ker:readouts}. That threshold may depend on
\(N\), the finite configuration and its margins. Only afterward take a
countable diagonal over increasing \(N\). The radius mesh is therefore
fixed before the discrete threshold is taken. This statement applies
to extremal-length tests in any deterministic units, as well as to the
other retained local readouts.
